\subsection{Relative safety and conserved helper theories}
\label{app:core-conservation-closures}

For a monotone Boolean function $f$ on a finite set $V$, let $S_f$ consist
of its marked cubes $(L,U,j)$, with $j\in U\setminus L$, on which
$f(L)=0$, $f(U)=1$, and $j$ is inessential. Write $R_f$ for the reflexive
transitive closure of inclusion and all their marked deletion edges,
and $P_f$ for the implication preorder of Boolean functions
having pure bichain formulas safe on $S_f$, in the sense of
\cref{app:core-grouped-cubes}. Endpoint constants are allowed.
Say that $f$ has \emph{relative conservation} if $P_f=R_f$: every upset of
the initial preorder already has a safe formula. The equivalence follows
from the bounded distributive lattice of safe Boolean functions.

Write $R_f^+$ for the preorder obtained by identifying all $f$-positive
vertices, and $R_f^-$ for the analogous identification of the negative
vertices. These identifications retain the original orientation from zero
to one.

Matroid spanning predicates have conserved theories with cofactor helper
bases in every Boolean realizer. For a finite
matroid $M$ of positive rank on $V$, let $h_M$ be the Horn closure with rules
\[
 L\longrightarrow j\qquad
 \bigl(L\text{ nonspanning},\quad j\in\operatorname{cl}_M(L)\bigr).
\]
A Boolean cofactor fixes selected variables consistently at endpoint constants.

\begin{theorem}[Matroid cofactor helper basis]
\label{core:thm:matroid-safe-prunings}
Let $f$ be the spanning predicate of $M$, and let $T$ be any pure Boolean
formula for $f$. Then $f$ has relative conservation, and
\[
 xR_fy\quad\Longleftrightarrow\quad x\subseteq h_M(y).
\]
The $S_f$-safe Boolean cofactors of $T$ separate every pair outside this
preorder. Hence every relative-safe Boolean function is a monotone
combination of these cofactors. The Boolean dual
$f^d(x)=1-f(V\setminus x)$ also has conservation,
with
\[
 xR_{f^d}y\quad\Longleftrightarrow\quad
 V\setminus y\subseteq h_M(V\setminus x),
\]
and the corresponding safe-cofactor assertion holds for every formula for $f^d$.
\end{theorem}
\begin{proof}
On a nonconstant cube $[L,U]$, $L$ is nonspanning and $U$ spans.
The marked $j$ is inessential exactly when
$j\in\operatorname{cl}_M(L)$, as in
\cref{core:thm:essential-cnf-matroid}. Every such cube can be enlarged to
$[L,V]$. Thus the initial edges are precisely the contextual reversals
of the displayed Horn rules. Reversing a Horn derivation and using
inclusion proves $xR_fy$ exactly when $x\subseteq h_M(y)$.

Suppose $j\in x\setminus h_M(y)$. Choose an inclusion-maximal
$h_M$-closed set $B\supseteq y$ avoiding $j$. This set contains a
matroid hyperplane $F$. If $B$ is nonspanning, Horn closure makes it
matroid-closed and keeps $j$ outside its closure. Extend a basis of $B$
together with $j$ to a matroid basis; the closure after removing $j$
is a hyperplane containing $B$ and avoiding $j$. It is $h_M$-closed,
so maximality makes it $B$ itself. If $B$ spans, choose an independent
$(\operatorname{rank}M-1)$-set $I\subseteq B$. Its closure is a
hyperplane, and the Horn rules give
$F=\operatorname{cl}_M(I)\subseteq B$.

The clause $q_B=\bigvee_{i\notin B}x_i$ separates $x,y$ and is $S_f$-safe.
Indeed a marked $j'$ in this clause could produce $\pv$ only if
$L\subseteq B$; inessentiality would then give
$j'\in\operatorname{cl}_M(L)\subseteq h_M(B)=B$, a contradiction.
It is an actual Boolean cofactor of $T$: fix $F$ to $c$ and
$B\setminus F$ to $b$. The remaining Boolean function is precisely
the disjunction on $V\setminus B$, since a hyperplane together with
any outside element spans. Every pure formula for this disjunction has
the same bichain semantics, by the uniform-threshold case of
\cref{core:cor:essential-unique-fiber}. Thus the cofactored term really
evaluates to $q_B$ on all bichain inputs.

All these safe cofactors respect $R_f$, and they separate every missing
pair. The full safe implication theory is therefore $R_f$. The bounded
lattice characterization gives the final generation assertion.
Swapping the endpoints and the two gates proves the dual statement.
\end{proof}

\begin{theorem}[Disjoint substitution]
\label{core:thm:conservation-substitution}
Let $F$ and $g_1,\ldots,g_m$ be nonconstant monotone Boolean functions,
with the $g_i$ on pairwise disjoint finite variable sets, and put
$f=F(g_1,\ldots,g_m)$. If all these functions have relative conservation,
then so does $f$.

More precisely, suppose $F$ is not a projection. Call $i$ an
\emph{AND-canalyzing} coordinate if $F|_{i=0}=0$, and
\emph{OR-canalyzing} if $F|_{i=1}=1$. Every other coordinate has both
cofactors nonconstant. The two canalyzing types cannot coexist. For conservation of $f$, it
suffices that $F$ and the $g_i$ at canalyzing coordinates have conservation,
while the remaining $D_{g_i}$ are term-realizable. Put
$T_i=R_{g_i}^+$ for an AND-canalyzing coordinate and
$T_i=R_{g_i}^-$ for an OR-canalyzing coordinate. Without any conservation
assumption, the initial preorder is
\[
 xR_fy\quad\Longleftrightarrow\quad
 g(x)R_Fg(y)\ \text{ and }\ x_iT_iy_i
 \text{ for every canalyzing }i,
 \qquad g(x)=(g_i(x_i))_i.
\]
For a projection $F=x_i$, it is $R_{g_i}$ on that block and the universal
preorder on all other blocks.
\end{theorem}
\begin{proof}
A cube in the product of the inner variable sets induces an outer cube:
each block contributes its constant value or a free outer coordinate.
This outer cube is nonconstant exactly when $f$ is nonconstant on the
original cube. If a marked inner coordinate changes $g_i$, its
$f$-inessentiality says that outer coordinate $i$ is inessential on this
outer cube. Thus every initial $f$-edge projects to an $R_F$-edge.

An AND-canalyzing coordinate is essential whenever it is free on a
nonconstant outer cube, and cannot be fixed zero there. Consequently
an initial deletion within its block either comes from a nonconstant
$g_i$-inessential cube or stays within the winning fiber. These are
$R_{g_i}^+$-edges. The OR case is dual. This proves the forward implication
in the displayed formula. It also shows that a canalyzing outer bit can
never decrease along $R_F$.

For the reverse implication, first observe that every $R_{g_i}$-edge
lifts with all other block vertices unchanged: enlarge those other
blocks to their full cubes. The inner nonconstant cube makes the
composite lower value zero and upper value one, and its marked coordinate
remains inessential. A constant-$c$ inner cube lifts in the same way
whenever $F|_{i=c}$ is nonconstant. Such cubes connect every pair of
vertices in the $c$-fiber: use their union when $c=1$ and their intersection
when $c=0$. Hence both fibers collapse for an interior coordinate,
whereas the winning fiber collapses for an AND-canalyzing coordinate
and the losing fiber for an OR-canalyzing coordinate. All these moves
are available independently of the other blocks.

An outer marked deletion can affect only an interior coordinate.
Choose any Boolean boundary edge of its $g_i$. With the other blocks
ranging over their required outer states, this gives a nonconstant
$f$-cube on which the marked inner boundary coordinate is inessential.
The two fiber collapses let it connect any prescribed inner vertices
of the two truth values. Outer inclusions lift similarly.

To coordinate these lifts with the canalyzing blocks, in the AND case
perform all prescribed $T_i$ paths first. Clamp each corresponding
outer bit that changes from zero to one to its final value one along
an $R_F$ path from $g(x)$ to $g(y)$. This remains a valid outer path:
every nonconstant semantic cube has that bit free or fixed one, so
its clamped vertices still belong to the original cube. Inclusion
steps are preserved as well. The other canalyzing bits remain constant,
and the interior steps lift as above. For OR-canalyzing coordinates,
clamp changing bits to their initial zero, lift the outer path first,
and perform the prescribed inner paths last. This proves the formula.
For a projection, ignored blocks are freely connected by the same
cube-lifting argument, giving the stated separate case.

It remains to separate every pair outside this preorder by a safe
formula. Use the assumed canonical realizers of the $g_i$, obtaining them from
conservation where needed: $g_i$ itself respects $R_{g_i}$, so conservation
gives an $S_{g_i}$-safe formula, and constant Boolean cubes have stable
endpoint values automatically.
An $S_F$-safe formula composed with these realizers is $S_f$-safe.
Indeed a marked inner source either is inessential within its block,
when the inner realizer removes it, or survives there and is inessential
on the induced nonconstant outer cube, when outer safety removes it.
Such substituted helpers separate any pair failing the outer comparison.

A pair failing a canalyzing comparison has a separating $T_i$-upset.
It is an $R_{g_i}$-upset constant on the saturated fiber, so conservation
of $g_i$ gives a safe formula for it. In the AND case, an $f$-nonconstant
cube sees that block either nonconstant or constantly one; its original
safety handles the first case and its constant winning-fiber value the
second. The OR case is dual. Marks in other blocks cannot appear in
this helper. These formulas separate all remaining missing pairs.
Since safe Boolean functions form a bounded lattice, their implication
preorder is therefore exactly $R_f$. The projection case follows by
lifting the inner safe helpers; constant outer functions have empty
budgets and are immediate.
\end{proof}

\begin{corollary}[Disjoint conjunction and disjunction]
\label{core:thm:conservation-disjoint}
Let $g,h$ be nonconstant monotone Boolean functions on disjoint finite
supports. Then
\[
 R_{g\land h}=R_g^+\times R_h^+,
 \qquad
 R_{g\lor h}=R_g^-\times R_h^-.
\]
If $g,h$ have relative conservation, so do $g\land h$ and $g\lor h$.
\end{corollary}
\begin{proof}
Apply \cref{core:thm:conservation-substitution} to the two Boolean gates.
Their initial preorders are inclusion, and their saturated factor
comparisons already imply the required comparison of truth values.
\end{proof}

\begin{corollary}[Read-once relative helpers]
\label{core:cor:readonce-relative-helpers}
Flatten the root of a nonconstant full-support read-once formula $f$ into
maximal opposite-root blocks $g_1,\ldots,g_m$, allowing individual variables.
For a projection take $m=1$ and $g_1=f$.
Then
\[
 xR_fy\quad\Longleftrightarrow\quad
 g_i(x)\leq g_i(y)\text{ for every }i.
\]
The Boolean functions of its relative-safe helpers are exactly
$H(g_1,\ldots,g_m)$ for arbitrary monotone Boolean $H$, including constants.
The original block formulas composed with any formula for $H$ give safe
realizers.
\end{corollary}
\begin{proof}
At an AND root each opposite-root OR block has one negative strongly
connected component, by the disjoint formula. Positive saturation therefore
reduces its preorder to the two-class truth preorder. The OR-root case
is dual. The disjoint supports make the block truth map surjective, and
the theorem supplies safety and all its upsets.
\end{proof}

\begin{corollary}[Read-once cofactor helper basis]
\label{core:cor:readonce-cofactor-helpers}
Let $f$ be a nonconstant full-support read-once Boolean function, and let
$T$ be any pure formula realizing $D_f$. Its $S_f$-safe Boolean cofactors
separate the initial preorder $R_f$ and generate all relative-safe Boolean
functions under monotone combination.
\end{corollary}
\begin{proof}
Use the maximal opposite-root blocks $g_i$ of
\cref{core:cor:readonce-relative-helpers}. Fix all variables in the other
blocks to the outer gate's neutral endpoint. Applying this same assignment
to $T$ gives a Boolean cofactor realizing $D_{g_i}$, since canonical
extensions commute with Boolean cofactors. The block functions $g_i$ are
$S_f$-safe and separate every pair outside $R_f$, by the preceding
corollary. Thus these actual cofactors of $T$ supply the claimed basis.
\end{proof}

\paragraph{The five minima of the path function.}
Initial completeness can also hold at every minimal realizer of a
noncanonical Boolean function. The four-variable path gives an exact
example; its classification of minimal realizers also proves the
cofactor failure below.

\begin{theorem}[Minimal path realizers and their conserved theories]
\label{core:thm:path-minimal-realizers}
On $V=\{A,B,C,D\}$ put $P=AB\vee BC\vee CD$. Its Boolean fiber has
exactly the following five pointwise-minimal operations:
\[
\begin{aligned}
 W_0&=B(A\vee C)\vee C(B\vee D),&
 W_1&=B(A\vee C)\vee CD,\\
 W_2&=AB\vee C(B\vee D),&
 W_3&=(A\vee C)(B\vee CD),\\
 W_4&=(B\vee D)(C\vee AB).
\end{aligned}
\]
Here products denote pure bichain meets. Every realizer of $P$ lies
above some $W_i$ in observer-union order. If $S_i$ consists of all
nonconstant $P$-inessential marked cubes on which $W_i$ is quiet, then
\[
 \preceq_{S_i}=R_{S_i}\qquad(0\leq i\leq4).
\]
All assertions allow observer union and endpoint constants.
\end{theorem}
\begin{proof}
Write $r_B=(D,BCD,B)$ and $r_C=(A,ABC,C)$. These are the only
nonconstant inessential cubes marked $B,C$, respectively: a free $B$
is pivotal whenever $A$ is available, or $C$ is available and $D$ is
not fixed true; reversal gives the $C$ case.
Put
\[
 p_A=(D,ACD,A),\quad u_A=(\varnothing,ACD,A),\quad
 q_D=(B,BCD,D),\quad \widehat q_D=(B,ABCD,D).
\]
Among nonconstant inessential marked cubes, the exact nonquiet rows are
\[
 \begin{array}{c|l}
 W_0&r_B,r_C\\
 W_1&r_B,q_D,\widehat q_D\\
 W_3&r_B,p_A,u_A.
 \end{array}
\]
For $W_0$, an $A$-source can survive only when $B$ is available and
$C$ is not fixed true, which witnesses actual $A$-essentiality by
setting $B=1,C=0$; $D$ is symmetric. Both $r_B,r_C$ give $\pv$.
For $W_1$ the same argument treats $A$, and $r_C$ is quiet.
An inessential $D$-source in its $CD$ summand survives exactly when
$B$ is fixed true and $C$ is free; nonconstancy permits $A$ to be
false or free, giving $q_D,\widehat q_D$. The row $r_B$ is nonquiet.
Path reversal $A\leftrightarrow D$, $B\leftrightarrow C$ exchanges
$W_1,W_2$ and $W_3,W_4$. Boolean duality followed by
$B\leftrightarrow C$ preserves $P$ and sends $W_1$ to $W_4$.
These symmetries give the remaining rows of the table and the exact
nonquiet sets of $W_2,W_4$.

No realizer is quiet on both $r_B,r_C$: their raw grouped cut gives
$BC\preceq AD$, contrary to $P(BC)=1$, $P(AD)=0$.
There is also a raw cut $CD\preceq BD$ on $\{p_A,r_C,q_D\}$.
Its upper routes to $p_A,q_D$ are inclusions, and the third is
\[
 CD\subseteq BCD\longrightarrow BC\subseteq ABC,
\]
where $q_D$ deletes $D$. Its lower routes are inclusions except
$A\subseteq AD\longrightarrow D\subseteq BD$, where $p_A$ deletes
$A$. The only removed coordinate is $C$, which is marked. Thus this
cut too is violated by $P$.

Let $T$ be any realizer. If it is nonquiet on both $r_B,r_C$, the table
gives $W_0\leq T$. Otherwise, by reversal assume $r_C$ quiet and $r_B$
nonquiet. The second cut forces $p_A$ or $q_D$ nonquiet. Enlarging a
cube preserves $\pv$, so the corresponding $u_A$ or $\widehat q_D$
is also nonquiet. The table now gives $W_3\leq T$ or $W_1\leq T$.
Reversal treats the other case. These comparisons use the one-source
order test of \cref{core:lem:essential-extension-pure}.
The five nonquiet sets are pairwise incomparable, so every $W_i$ is
minimal and the list is complete.

It remains to prove initial completeness. Put
$L=B(A\vee C)$, $M=C(B\vee D)$. On $S_0$, the initial rules delete
$A$ exactly when $B=0$ or $C=1$, and delete $D$ exactly when $C=0$
or $B=1$; neither $B$ nor $C$ can decrease. The six equivalence classes
have representatives and safe flag vectors
\[
 \begin{array}{c|cccccc}
 \text{representative}&\varnothing&B&AB&C&CD&BC\\\hline
 (B,C,L,M)&0000&1000&1010&0100&0101&1111.
 \end{array}
\]
Indeed at fixed $B,C$ the displayed rules clear $A,D$ except for $A$
on the $B=1,C=0$ branch and $D$ on the $B=0,C=1$ branch.
The initial order is precisely componentwise order on these vectors:
the two chains are $\varnothing<B<AB<BC$ and
$\varnothing<C<CD<BC$. The flags exclude further comparisons.
They are all safe: $B,C$ are unmarked on $S_0$, and $L,M$ are root
children of a quiet join on nonconstant $P$-cubes. Their monotone
combinations therefore realize every upset of $R_{S_0}$.

On $S_1$ the initial rules instead delete $A$ when $B=0$ or $C=1$,
delete $D$ when $C=0$, and delete $C$ when $A=1,D=0$; $B$ cannot
decrease. Their classes are
\[
 \begin{array}{c|ccccc}
 \text{representative}&\varnothing&CD&B&AB&BCD\\\hline
 (B,L,CD)&000&001&100&110&111.
 \end{array}
\]
The rules identify $AB$ and $BC$ when $D=0$, and clear a lone $C$
with $B=D=0$ by $C\to AC\to A\to\varnothing$.
The remaining reductions just clear $A$ or $D$ by the stated rules.
This proves that each flag fiber is one class. Inclusion and these
reductions give exactly componentwise order on the five vectors.
Here $B$ is unmarked, while $L,CD$ are the two root children of $W_1$,
so all three flags are safe. Again they generate every initial upset.
The bounded lattice characterization proves
$\preceq_{S_0}=R_{S_0}$ and $\preceq_{S_1}=R_{S_1}$.
Reversal and Boolean duality transfer the second equality to the
other three minima. Purification covers observer union.
\end{proof}

\begin{theorem}[Conservation through four variables]
\label{core:thm:four-variable-conservation}
Let $f$ be monotone Boolean on a variable set $V$ with $|V|\leq4$.
Unless $f$ is
$AB\vee BC\vee CD$ up to variable permutation, it has relative
conservation and a realizable canonical extension. The path function has
no realizable canonical extension.
For every pointwise-minimal realizer $T$ of $f$ on $V$, let $S_T$ be
its quiet nonconstant $f$-inessential marked cubes. Then
\[
 \preceq_{S_T}=R_{S_T}.
\]
The realizer and helpers use variables in $V$; observer union and
endpoint constants are allowed.
\end{theorem}
\begin{proof}
Constants are immediate. For a nonconstant function, its marked full cube
permits deletion of any unused coordinate everywhere. The initial preorder
and its upsets therefore ignore that coordinate, reducing conservation
to the essential variables. On at most three of these the function is
read-once or a uniform threshold:
a singleton minterm splits off its variable, and otherwise the only
full-support cases are a conjunction, a two-edge star, or majority.
All these functions have conservation.
On four variables, a singleton minterm or singleton prime clause reduces
to a disjoint OR or AND with a function on at most three variables.
Use \cref{core:thm:conservation-disjoint}.

Otherwise minterms have size two or three. If there is no pair, avoiding
a common variable requires all four triples, giving the three-of-four
threshold. Let $G$ be the graph of pair minterms. If a triple minterm
occurs, it contains no edge of $G$, so every edge meets its omitted
vertex: $G$ is a star. With three star edges the only possible extra
triple gives
\[
 ABC\vee Z(A\vee B\vee C),
\]
an apex extension. With two edges, the extra-triple case is
$Z(A\vee B)\vee ABC$, whose Boolean dual is the apex
$AB\vee Z(A\vee B\vee C)$. With one edge $ZA$, a single extra triple
has a common variable and was already treated; both possible triples give
$ZA\vee BC(A\vee Z)=\operatorname{Maj}_3(Z,A,BC)$.
Thus \cref{core:thm:conservation-apex,core:thm:conservation-substitution}
and Boolean duality prove conservation in every mixed case.

It remains to treat quadratic functions. A four-vertex cograph has an
isolated vertex, a universal vertex, or is $2K_2$ or $C_4$.
An isolated vertex reduces arity; a universal vertex gives an apex
over a quadratic function on at most three variables. The last two
graphs give $AB\vee CD$ and $(A\vee B)(C\vee D)$, disjoint compositions.
The only remaining graph is the path, whose canonical extension is
excluded by \cref{core:thm:cograph-extension}.

In the conserved cases the canonical extension is the least member of
the Boolean fiber, hence its only minimal member, with quiet budget
$S_f$. For the path, \cref{core:thm:path-minimal-realizers} proves the
assertion for every minimal member. Purification covers observer union.
\end{proof}

Canonicity passes to Boolean cofactors, but pointwise minimality need not.
The example above attains the smallest possible number of essential inputs.

\begin{corollary}[Sharp failure of minimality under Boolean cofactors]
\label{core:cor:minimal-cofactor-obstruction}
The minimal realizer $W_0$ has a nonminimal cofactor at $A=c$:
\[
 W_0|_{A=c}=B\vee\bigl(C\wedge(B\vee D)\bigr)
    > B\vee(C\wedge D)
\]
in pointwise observer-union order.
Four essential Boolean variables are necessary for such a failure,
even allowing observer union and endpoint constants.
\end{corollary}
\begin{proof}
The read-once word $B\vee(C\wedge D)$ is the canonical lower envelope
of this cofactor. At $(B,C,D)=(\nv,\pv,b)$ it returns $\nv$, whereas
the cofactor returns $\pv$, proving strictness.

With at most three essential variables the preceding classification gives
a canonical realizer that ignores all unused inputs. It is the global
bottom and unique minimal member of its Boolean fiber. Boolean cofactors
of a canonical extension are canonical by the cube definition, proving
the lower bound.
\end{proof}

\paragraph{Lattice generation need not be a single pruning.}
An occurrence-constant pruning independently retains each variable leaf
or replaces it by an endpoint. Such prunings need not be closed under
Boolean meet with the original function, even when every pruning is safe.

\begin{remark}[Single-pruning saturation can fail]
\label{core:rem:pruning-saturation}
The canonical threshold word
\[
 T=(a\vee b\vee c)\wedge(d\vee ab\vee ac\vee bc)
\]
computes $f=\operatorname{Thr}_2(\{a,b,c,d\})$ and has a pruning with
Boolean function $p=a\vee bc$. Nevertheless no pruning of $T$ computes
\[
 f\wedge p=ab\vee ac\vee ad\vee bc=:h.
\]
Here $S_f$ is empty, so all prunings are relative-safe. The target $h$
also has a canonical extension, by \cref{core:thm:cograph-extension}.
\end{remark}
\begin{proof}
At $d=0$, the right factor of $T$ requires two of $a,b,c$; at $d=1$,
the left requires one. Thus $T$ is the indicated uniform threshold and
has its unique canonical semantics. To obtain $p$, prune the left
factor to $1$, remove $d$ and $ab$, prune $ac$ to $a$, and retain $bc$.

Any pruning is $\ell\wedge r$, where $\ell$ is constant or a disjunction
of a subset of $a,b,c$, and $r$ is constant or a disjunction of monomials
from $a,b,c,d,ab,ac,bc$. If it computed $h$, then $h(ad)=1$ would give
$\ell(a)=1$. Since $h(a)=0$, it follows that $r(a)=0$.
Thus $r(ad)=1$ forces a standalone $d$ monomial. Now $h(d)=0$ excludes
$\ell=1$, while $\ell(ad)=1$ forces $a$ into its support and
$\ell(bc)=1$ forces $b$ or $c$ into that support. Together with the
$d$ monomial of $r$, this makes $bd$ or $cd$ true, a contradiction.

The graph of $h$ is the paw, a cograph. Its four monomials $ab,ac,ad,bc$
are separately obtainable as prunings of $T$, so their disjunction does
belong to the bounded lattice generated by safe prunings. The obstruction
concerns a single pruning. Swapping meet with join and the endpoints
gives the dual failure of join saturation for a three-of-four threshold.
\end{proof}

\begin{corollary}[Majority wrapper]
\label{core:cor:conservation-majority-wrapper}
Let $g$ be a nonconstant monotone Boolean function and $u,v$ fresh variables.
Put $W=\operatorname{Maj}_3(g,u,v)$. Then
\[
 (x,u,v)R_W(y,u',v')\quad\Longleftrightarrow\quad
 g(x)\leq g(y),\quad u\leq u',\quad v\leq v'.
\]
Moreover $W$ has relative conservation if and only if $D_g$ is
term-realizable. Consequently, conservation is preserved by all Boolean
cofactors if and only if every term-realizable canonical extension has
relative conservation.
\end{corollary}
\begin{proof}
Every free variable of $\operatorname{Maj}_3$ is essential on each of its
nonconstant cubes. Its semantic budget is therefore empty, its initial
preorder is inclusion, and it has conservation. Both cofactors at every
coordinate are nonconstant, so \cref{core:thm:conservation-substitution}
gives the displayed preorder and proves conservation of $W$ from a
canonical realizer of $g$. Conversely, conservation makes $W$ itself
safely realizable, hence makes $D_W$ realizable. Fixing $u=0,v=1$ in this realizer gives $D_g$, since the marked-cube
semantics restricts to that of $g$.

If conservation is hereditary under cofactors, apply it to this wrapper
to obtain conservation for any canonically realizable $g$. Conversely,
every conservative function has a canonical realizer, and canonical
realizability survives Boolean cofactors; the asserted universal property
then gives conservation of each cofactor. Constant functions have empty
budgets and cause no exception.
\end{proof}

There is also a closure with overlap across every old variable. Let
$p$ be monotone on finite nonempty $V$ with $p(\{i\})=0$ for every $i$,
and put $s=\bigvee_{i\in V}x_i$. The zero function is permitted.

\begin{theorem}[Apex extension]
\label{core:thm:conservation-apex}
If $p$ has relative conservation, then
\[
 F(z,x)=p(x)\lor\bigl(z\land s(x)\bigr)
\]
has relative conservation. The literal composition
$D_p\lor(z\land s)$ is its canonical essential-source extension.
\end{theorem}
\begin{proof}
On a cube $[L,U]$ with $L\ne\varnothing$, $s$ is constantly one and
the composition reduces to $D_p\lor z$. If $L=\varnothing$, every free
old coordinate is essential when $z$ is free or fixed to one: set $z=1$
and use $s$. On a nonconstant cube with $z$ free it is essential as well,
since an available singleton has $p=0,s=1$. With $z=0$ use $D_p$.
These cases prove canonicality.

Define two budgets on $V$:
\[
\begin{aligned}
 B_1&=\{(L,U,j):L\ne\varnothing,\ p(L)=0,
                   \ j\text{ inessential on }[L,U]\},\\
 B_0&=S_p\cup\{(L,U,j):L\ne\varnothing,\ p(U)=0\}.
\end{aligned}
\]
Let their initial preorders be $R_1\subseteq R_0$. No nonconstant
$F$-cube has an inessential $z$. Cubes with $z=0$ give $S_p$; cubes
with $z=1$ give no marks; cubes with $z$ free give precisely $B_1$.
Consequently the two faces of $R_F$ carry $R_0,R_1$, respectively, and
\[
 (0,x)R_F(1,y)\Longleftrightarrow xR_0y,
 \qquad (1,x)\not R_F(0,y).
\]

We prove conservation for $B_1$. Join distinct $i,j\in V$ by an edge
when $p(\{i,j\})=0$, and let $\mathcal C$ be the set of components.
Every nonempty zero set lies in one component. All nonempty zero sets
within a component are $R_1$-equivalent: constant-zero cubes with
nonempty lower endpoints identify a zero set with its singletons, and
zero pairs connect the singletons along graph paths. Different components
cannot connect through the zero region, and no positive vertex can
return to zero. Empty is alone.

Every $R_p$ comparison between positive vertices survives in $R_1$.
For a marked deletion with empty lower endpoint and positive target $A$,
restrict the lower endpoint to any $\ell\in A$. It remains nonconstant,
since $p(\{\ell\})=0$, and retains the deletion. Positive paths never
leave the positive region.

Let $\phi$ be an $R_1$-upset. If $\phi(\varnothing)=1$, use the constant
formula. Otherwise let $b_C$ be its common value on the nonempty zero
sets in $C\in\mathcal C$, and let $\psi$ equal $\phi$ on positive vertices and zero
on all zero vertices. The preceding positive-path observation makes
$\psi$ an $R_p$-upset. A $p$-relative-safe realizer of $\psi$ is also
$B_1$-safe, since it is constantly zero on every added zero cube.

Each $s_C=\bigvee_{j\in C}x_j$ is $B_1$-safe. Otherwise a marked
$j\in C$ would have $L\cap C=\varnothing$. Choose $\ell\in L$ outside
$C$. Then $p(\{\ell,j\})=1$, while $p(L)=0$, making $j$ essential
already at the lower corner, a contradiction. Finally
\[
 \phi=\psi\lor\bigvee_{C\in\mathcal C:\,b_C=1}s_C.
\]
This is exact on zero vertices; on positive vertices monotonicity from
the activated singleton values makes every added $s_C$ lie below $\phi$.
Thus $B_1$ is conserved.

An $R_0$-upset has a $p$-relative-safe realizer and is constant on every
added zero cube, so $B_0$ is conserved too. An $R_F$-upset is a pair
$u_0\leq u_1$ of $R_0$- and $R_1$-upsets. Choose safe realizers
$T_0,T_1$. Then $T_0\lor(z\land T_1)$ realizes the pair and is safe on
all $F$-rows: on a $z=0$ row use $T_0$; on a free-$z$ marked row both
old formulas are quiet and $z$ is an unmarked neutral input. This proves
conservation.
\end{proof}

\begin{proposition}[Threshold apex reduction]
\label{core:prop:conservation-threshold-apex-reduction}
Let $1\leq k\leq|V|$ and suppose $p(A)=0$ whenever $|A|\leq k$.
If $D_p$ is term-realizable, then
\[
 D_p\lor\bigl(z\land D_{\operatorname{Thr}_k(V)}\bigr)
   =D_{p\lor z\operatorname{Thr}_k(V)}.
\]
Its initial preorder has the same triangular form, with
\[
 B_k=\{(L,U,j):|L|\geq k,\ p(L)=0,
                    \ j\text{ inessential on }[L,U]\}
\]
on the $z=1$ face, and $R_p$ augmented by the constant-zero cubes of
$B_k$ on the $z=0$ face.
\end{proposition}
\begin{proof}
With $z=0$ use $D_p$. If $|L|\geq k$, the threshold is one and the
operation is $p\lor z$.
If $|L|<k\leq|U|$, every free old coordinate is essential by setting
$z=1$. A $k$-set between $L,U$ has $p=0$ and threshold one, so $z$ is
essential too when free. If $|U|<k$, the whole function is zero. These cases give
canonicality and the asserted marks.
\end{proof}

\begin{corollary}[Connected zero faces]
\label{core:cor:conservation-connected-zero-faces}
Under the hypotheses of \cref{core:prop:conservation-threshold-apex-reduction},
suppose that $p$ has relative conservation. Form the graph on the
$k$-subsets of $V$, joining distinct $A,B$ when $p(A\cup B)=0$.
If this graph is connected, then $B_k$ and
$p\lor z\operatorname{Thr}_k(V)$ have relative conservation.
In particular, this holds if $p(A)=0$ whenever $|A|\leq k+1$.
\end{corollary}
\begin{proof}
Write $R_k$ for the initial preorder of $B_k$. Its zero vertices of size
at least $k$ form one strongly connected component $C$. A constant-zero
cube with lower endpoint of size $k$ identifies each such vertex with
its $k$-subsets; each graph edge identifies its endpoints through their
zero union. Every budget edge preserves $p$.

Every $R_p$ comparison between positive vertices survives in $R_k$.
For a marked deletion with positive target $A$ and lower endpoint $L$
of size less than $k$, choose $L\subseteq K\subseteq A$ with $|K|=k$.
The restricted cube remains nonconstant because $p(K)=0$, and retains
the deletion. Paths starting at positive vertices remain positive.

Let $\phi$ be an $R_k$-upset. If $\phi(C)=0$, then $\phi$ vanishes at
every zero vertex: a vertex of size less than $k$ extends to a zero
$k$-set. The positive-comparison observation therefore makes $\phi$
an $R_p$-upset. A $p$-relative-safe realizer is also $B_k$-safe,
since it is stably zero on every added constant-zero cube.

If $\phi(C)=1$, monotonicity gives
$\phi\geq\operatorname{Thr}_k(V)$. Choose monotone formulas $M,Q$
for $\phi,\operatorname{Thr}_k(V)$. Then $M\lor Q$ still computes
$\phi$ and is $B_k$-safe: on every budget cube $Q$ is stably one
because $|L|\geq k$. Thus every $R_k$-upset has a safe formula.
The other face preorder is conserved by the constant-zero-cube
argument, and the triangular synthesis proves conservation of the
threshold extension.

Finally, if all $(k+1)$-sets are zero, exchanging one element of a
$k$-set at a time gives paths in the graph between any two vertices.
The one-vertex case is included.
\end{proof}

\begin{theorem}[A required input]
\label{core:thm:conservation-required-input}
Let $1\leq k\leq|V|$, let $G$ be nonconstant and monotone on $V$, and
suppose $G(A)=0$ whenever $|A|<k$. Let $t,z$ be new inputs and put
$p=t\land G$. If $p$ has relative conservation, then so do $B_k$ and
\[
 p\lor z\operatorname{Thr}_k(V\cup\{t\}).
\]
Here $B_k$ is the budget from
\cref{core:prop:conservation-threshold-apex-reduction}, on $V\cup\{t\}$.
\end{theorem}
\begin{proof}
All $k$-subsets of $V$ are zero vertices of $p$ and lie in one zero-face
component $C$: their unions omit $t$. A remaining $k$-face is
$\{t\}\cup B$ with $|B|=k-1$. If a zero $k$-set of $G$ contains $B$,
this face connects to $C$ through that set. Otherwise every $k$-set
containing $B$ is $G$-positive, and $\{t\}\cup B$ is isolated: its union
with any other $k$-face is positive. These are all the components.
Every zero vertex of size at least $k$ is identified with its $k$-faces
by constant-zero cubes. Also write $C$ for the resulting component of
the initial preorder $R_k$ of $B_k$.

Let $D$ be the upward closure of the faces in $C$, viewed as a monotone
Boolean function. It is stably one on every $B_k$ cube. Indeed, a zero
lower endpoint $L$ of size at least $k$ with $D(L)=0$ must be one of the
isolated faces $\{t\}\cup B$. At that lower endpoint every free old
coordinate is essential, since adjoining it makes $G$ positive. Hence
such a cube has no inessential mark.

Let $\phi$ be an $R_k$-upset. If $\phi(C)=1$, then $\phi\geq D$.
The join of any monotone
formulas for $\phi,D$ still computes $\phi$ and is safe on $B_k$, because
the second formula is stably one there.

Suppose instead that $\phi(C)=0$. It vanishes on every vertex omitting
$t$: small old sets extend to old $k$-sets. The $p$-positive vertices
form one $R_p$-component by \cref{core:thm:conservation-disjoint}.
Their comparisons survive in $R_k$ by the positive-path argument in
\cref{core:cor:conservation-connected-zero-faces}. Thus either $\phi=0$,
or it is one at every $p$-positive vertex.

In the latter case, any active zero vertex $\{t\}\cup S$ has $|S|<k$.
Moreover every old $k$-set containing $S$ is $G$-positive; otherwise its
union with $\{t\}$ would be a zero vertex in $C$ above this active vertex.
For such an $S$ define
\[
 \chi_S=G\lor\bigwedge_{i\in S}x_i.
\]
This is an $R_G$-upset. A violating marked deletion $A\cup\{j\}\to A$
would have both endpoints $G$-zero, with $j\in S\subseteq A\cup\{j\}$.
Universality of $S$ forces $|A\cup\{j\}|<k$. If $[L,U]$ witnesses this
deletion, choose a $k$-set $K\subseteq U$ containing $L\cup S$.
Then $G(K)=1$, whereas $G(K\setminus\{j\})=0$; both sets lie in the
cube. This contradicts the inessentiality of $j$.
Since $\chi_S$ is one on every $G$-positive vertex, it also respects
$R_G^+$.

Consequently each $t\land\chi_S$ is an $R_p$-upset, and so is
\[
 \phi=p\lor
   \bigvee_{\substack{G(S)=0\\\phi(\{t\}\cup S)=1}}
       \left(t\land\bigwedge_{i\in S}x_i\right).
\]
Conservation of $p$ supplies a formula safe on $S_p$. It is also safe
on every added constant-zero cube of $B_k$: an isolated zero $k$-face
has no larger zero superset, so every such cube lies entirely in $C$,
where $\phi$ is zero. Thus every $R_k$-upset has a $B_k$-safe formula.

The other face preorder is conserved by the constant-zero-cube argument.
The triangular synthesis from the threshold-apex reduction now proves
conservation of the displayed extension.
\end{proof}
